\section{Label Budgets, Joint Naming and Verification}
\label{sec:quality}

\subsection{Annotation resources and joint naming}

$L_m$ holds the m labelled development examples of Section~\ref{sec:problem};
test queries are inductive and outside training graphs.

Fourteen public learners include regional majority, three assignment variants
(differing in how empty regions are treated), class centroids, five ridge
heads and four graph propagation configurations \cite{zhou2003consistency}
with two strengths, with/without class-mass normalization. Ridge and graph
learners use exact LOO predictions; a fixed public-first order resolves ties.
Web selection uses observed-class balanced accuracy.

With observed regional label counts $L[k,c]$, projected private matrix A and
public region sizes $w_k$, joint naming solves
\begin{equation}
\hat\pi_\lambda=\arg\max_{\pi\in\mathfrak S_C}
\sum_k\left\{L[k,\pi(k)]+\lambda\frac{w_k}{\bar w}A[k,\pi(k)]\right\}.
\label{eq:fusion}
\end{equation}
Here $K=C$, $\bar w=\sum_k w_k/K$ and
$\lambda\in\{.01,.03,.1,.3,1,3,10,30\}$. LOO selects among the fourteen
public learners, joint assignments and fixed private centered/coverage rules.
All candidates reuse one protected matrix, so selection adds no private
query. Exact LOO is efficient for the assignment candidates.

\begin{valueprop}[LOO on matched edges]
\label{prop:loo}
For $S=L+\lambda\operatorname{diag}(w/\bar w)A$, let $\pi$ maximize assignment
weight. Removing an example on an unmatched edge $(k,c)$, $c\ne\pi(k)$,
cannot make that edge correct in the reduced problem. Thus LOO correctness
for all m examples needs at most $C+1$ assignment solves per weight.
\end{valueprop}
Every assignment using an unmatched removed edge loses one unit while the
optimum loses none; examples on one matched edge share one refit.

The same joint decoder receives a protected head's mean proxy-region softmax
or point scores averaged inside each region before projection. Native head
prediction is an additional fixed head-arm candidate.

\subsection{Public evidence and publication targets}

The public augmentation rule requires a positive paired LOO lower gain,
all-class proxy predictions and a positive Wilson-form recall bound per class
(heuristics, since prediction and selection are dependent), else it falls
back to the public model. A missing labelled class forces rejection, so
verification labels also strengthen the public competitor; private
verification changes the source of evidence and charges its privacy.

\subsection{Private verification within one total budget}
\label{sec:private-verification}

After protected naming fixes candidate f and public comparator b, client i
computes held-out gain $d_i=u_i(f)-u_i(b)\in[-1,1]$ and a normalized
$C\times4$ histogram of true class and joint correctness of f and b; both
use distributed discrete-Gaussian shares.

The demonstration fixes $\rho$ fractions $1/2,1/4,1/4$ for naming, gain and
class statistics. Conditional composition in Theorem~\ref{thm:run} charges
all three releases to the same total $(\varepsilon,10^{-5})$ budget.
With released mean $\tilde r$ and gain-channel mean-noise proxy $v_g$, set
\begin{equation}
\underline G=2\tilde r-1
-2v_g\sqrt{2\log(1/\alpha_g)}-s^{-1}.
\label{eq:private-gain-bound}
\end{equation}
Discrete-Gaussian sub-Gaussian tails \cite{canonne2020discrete} imply
$\Pr(\underline G>G)\le\alpha_g$, where $G=n^{-1}\sum_i d_i$.
Direct histogram numerator/denominator bands yield simultaneous class-recall lower
bounds (Proposition~\ref{prop:private-validation} below). We fix
$\alpha_g=\alpha_h=.025$: gain-only acceptance uses $\underline G>0$;
class-checked acceptance also requires every recall lower bound positive. A
\emph{false accept} is an accepted case whose clean held-out statistics miss
the target ($G\le0$, or some class recall zero or below the threshold).
Rejection returns the public comparator without certifying its class
recall. Coverage concerns validation noise and the finite held-out cohort,
not unseen clients; publishing or rejecting is post-processing.

\paragraph{Bounded validation reports}
In the cost ladder each client has 120 held-out records. For paired gain, encode
$q_i=\operatorname{round}(s(1+d_i)/2)\in[0,s]$. Its integer sensitivity
is s, and converting the rounded mean to gain adds at most $1/s$ error.
The class report $h_i\in\Delta^{4C-1}$ records true class and joint binary
correctness $(a,b)$ of candidate/comparator. Largest-remainder rounding
gives $Q_i\ge0$, $\sum Q_i=s$, with coordinate error below $1/s$.
Its integer squared sensitivity is exactly $2s^2$. Both reports have fixed
dimension (1 and $4C$), independent of the number of held-out records.

\begin{valueprop}[Conditional finite-heldout coverage]
\label{prop:private-validation}
Assume the specified reports and independent noise shares. Fix the naming
transcript, public labels and actual held-out cohort.
Equation~\eqref{eq:private-gain-bound} has one-sided coverage at least
$1-\alpha_g$. Let $z=\sqrt{2\log(4C/\alpha_h)}$ and define
\begin{equation}
a_N=\sqrt2v_h z+2/s,\qquad a_D=2v_h z+4/s.
\label{eq:hist-radius}
\end{equation}
For released normalized histogram $\tilde h$, take
$D_c^U=\min(1,[\sum_{a,b}\tilde h_{c,ab}+a_D]_+)$ and
$N_c^L=[\tilde h_{c,10}+\tilde h_{c,11}-a_N]_+$.
The clipped ratio $N_c^L/D_c^U$ is simultaneously a recall lower bound for
all classes with probability at least $1-\alpha_h$, taking zero if $D_c^U=0$.
Joint gain/class coverage is at least $1-\alpha_g-\alpha_h$, apart from
the stated wrap term.
\end{valueprop}
\begin{proof}[Proof sketch]
Each exact share is sub-Gaussian with proxy $b$ (Lemma~16 of
Canonne et al.~\cite{canonne2020discrete}), so the mean has proxy $v$; a
correct-mass sum has proxy $\sqrt2v_h$ and a class mass $2v_h$. A union over
the $2C$ fixed linear targets with rounding errors gives
Eq.~\eqref{eq:hist-radius}, and monotone division plus a union with the gain
event completes the claim.
\end{proof}

This is noise coverage for a finite cohort, not a posterior probability
given acceptance. Whole-client adjacency affects training and held-out
records, so all releases compose sequentially.

\paragraph{Why targets differ in cost}
The gain target resolves one scalar with multiplier
$\sqrt{2\log(1/\alpha_g)}$; classwise targets resolve $2C$ simultaneous
bands with the larger $z$, and each must separate from zero a class mass of
order $1/C$. Section~\ref{sec:eval:verification} measures the resulting
cost ladder.

Higher K can raise purity while increasing the noise in
Eq.~\eqref{eq:design-budget}; natural writers test this.
